\section{Introduction}

The quality of honey is assessed not only by instrumental analysis but also
by the senses of sight, smell and taste. Organoleptic assessment is
inexpensive and quick to perform, and its result has practical consequences,
because it may decide whether a batch is admitted to the market. At the same
time, it is the most difficult part of the assessment to standardize, since
it relies on human judgment rather than on an instrument reading. The legal
framework for this assessment is set by European Union law, national
regulations and international standards.

The commercial quality of honey placed on the market in the European Union is
defined by Directive 2001/110/EC \cite{Dyrektywa2001110}. The directive is
transposed into the law of the Member States, and outside the Union the
reference point is the Codex Alimentarius standard CXS 12-1981
\cite{CodexHoney}. Besides instrumentally determined parameters, both
documents contain organoleptic requirements. Honey must not have any foreign
taste or odor, and a floral (varietal) name may be used only if the product
has the organoleptic, physicochemical and microscopic characteristics of the
declared source.

Polish regulations impose a broader requirement. The regulation on the
commercial quality of honey \cite{RozporzadzenieMRiRW2003} considers honey
compliant if it ``has an odor and taste typical of the given honey variety''.
A finding of atypical taste, however, requires a comparison with a reference
for the variety, which none of these acts defines. Descriptive sheets of the
main European unifloral honeys are available in the literature
\cite{PersanoOddo2004}, and a methodology of organoleptic assessment has been
published within the International Honey Commission (IHC) \cite{Piana2004}.
These descriptions do not, however, have the status of a legal reference. At
the same time, Article~5 of the directive \cite{Dyrektywa2001110} prohibits
Member States from adopting national provisions on honey not provided for by
the directive.

In Poland, the method of organoleptic assessment is laid down in the
regulation on analytical methods for honey \cite{RozporzadzenieMRiRW2009}.
The regulation only states that color, consistency, taste and odor are
examined with the senses of sight, smell and taste, and that the results are
compiled in a table. It specifies neither the number nor the qualifications of
the assessors, nor a reference material, nor a criterion for classifying a
taste as atypical. Directive (EU) 2024/1438, applicable from 14~June 2026
\cite{Dyrektywa20241438}, provides for the adoption of harmonized methods for
detecting honey adulteration, but not of a method for organoleptic
assessment.

How the assessment is performed and how its result is interpreted therefore
remains at the discretion of the laboratory. As a consequence, an organoleptic
verdict may by itself determine the non-compliance of a batch, even when all
other determinations raise no objections.

% A criterion with legal and administrative consequences is expected to be
% objective, reproducible and proportionate. In official controls this
% requirement has a legal basis: laboratory methods must comply with Union
% rules or, in their absence, with internationally recognized protocols, with
% national rules, or be validated in inter- or intralaboratory studies
% \cite[Art.~34]{Rozporzadzenie2017625}. In the 2015 coordinated control plan
% for honey of the European Commission, the organoleptic test was included in
% the first stage of sample characterization, together with electrical
% conductivity, diastase activity and pollen analysis, as a test indicating
% the direction of further, more specific determinations
% \cite{KE2015PlanKontroli}. The measurability of a characteristic does not
% yet establish its suitability as a decision criterion: this was shown for
% the microbiological parameters of honey, where the count of osmophilic
% yeasts alone, without signs of active fermentation, is not a sufficient
% basis for rejecting the product \cite{Gajda2026}. The same question can be
% asked about the assessment of taste and odor typicality.

The psychophysical literature points to two independent sources of
variability in organoleptic assessment. The first is the interindividual,
partly heritable variability in the perception of basic tastes: the
recognition threshold for sourness is about 53\% heritable \cite{wise2007twin},
and polymorphisms of the \textit{GNAT3} gene explain about 13\% of the
variance in the perception of sucrose sweetness \cite{fushan2010sweet}. The
same sample therefore evokes sensations of genuinely different intensity in
different people, and perceptual experiences cannot be compared directly across
individuals \cite{bartoshuk2000comparing}. The second source is the
non-comparability of descriptive scales: verbal labels (e.g., ``very strong'',
``typical'') denote different absolute intensities for different people,
which makes category scales an inappropriate tool for comparisons between
individuals and groups \cite{bartoshuk2003labeled}.

The aim of the study was to test empirically whether the assessment of honey
taste and odor typicality is reproducible enough to serve as a stand-alone
criterion of commercial quality. To this end, the results of a blind consumer
panel were compared with the organoleptic assessments of the same samples
carried out independently by laboratories. The analysis focused on the ability
of the blind consumer panel to discriminate between samples in terms of taste
typicality, on the agreement of the assessments between laboratories, and on
the extent to which consumers and laboratories identified the same samples as
not typical.
